%============================================================================
% prelude.tex
%   - titlepage
%   - table of contents, list of tables and list of figures
%   - abstract
%============================================================================
\clearpage%
\pagenumbering{roman}  % This makes the page numbers Roman (i, ii, etc)
%--------------------------------------------------------------------%
% TITLE PAGE
%   - define \title{} \author{} \date{}

\title{ Analysis using Belief Propagation Techniques }
\author{Chitra Suresh}

\date{December 2014}

\department{Department of Electronics Engineering}

%    - Set the guide's name
\setguide{Prof. Dr.Kushal.R.Tuckley}
%    - Set the coguide's name (if you have one)
%%\setcoguide{Prof Amitabha Sanyal}

%   - once the above are defined, use \maketitle to generate the titlepage
\maketitle

%--------------------------------------------------------------------
\cleardoublepage \nonumber
%\begin{center}
%\includegraphics{raitlogo.eps}\\
%Ramrao Adik Education Society's\\
%\large {\textbf{Ramrao Adik Institute of Technology}}\\
%\small(Affiliated to the University of Mumbai)\\
%Dr. D. Y. Patil Vidyanagar,Sector 7, Nerul, Navi Mumbai 400
%706.\\\vspace{0.3in} \Large CERTIFICATE
%\end{center}
%\vspace{0.1in}
%\begin{center}
%\large This to certify that Dissertation Seminar-I entitled\\
%\textbf{\lq\lq On Scheduling in WiMAX\rq\rq} \\submitted by
%\begin{center}
% w (09)\\
%\end{center}
%is approved for the degree of \\ \textbf {Master of Engineering}\\ in\\ \textbf {Electronics Engineering}.\\
%
% %\vspace{0.5in}
%%-----------------------------------------------------------------
%%\begin{tabular}{ccc}
%%      \rule{6cm}{1sp}                &\rule{10mm}{0pt}& \rule{6cm}{1sp}
%%      \\
%%       Guide             &&Head of Department \\\\\
%%      \rule{6cm}{1sp}                && \rule{6cm}{1sp} \\
%%      Examiner 1                 && Examiner 2 \\\\\\
%%      Principal: \rule{4cm}{1sp} && \\ && \\
%%       &&
%%    \end{tabular}
%\vspace{1in}
%\begin{tabular}{ccc}
%      \rule{5cm}{1sp}                &\rule{10mm}{0pt}& \rule{5cm}{1sp}
%      \\\vspace{0.5in}
%      Examiner 1                 && Examiner 2 \\
%      \rule{5cm}{1sp}                && \rule{5cm}{1sp} \\ \vspace{0.5in}
%       Guide             && Project Coordinator \\
%      \rule{5cm}{1sp}                && \rule{5cm}{1sp} \\
%      Head of Department                && Principal \\\\
%        %&\rule{5cm}{1sp}& \\
%%       &Principal&
%    \end{tabular}
%\end{center}
%-----------------------------------------------------------------
\cleardoublepage
%--------------------------------------------------------------------%


% ABSTRACT
\begin{abstract}
%\input {abstract}
Probabilistic models have received considerable attention,particularly as they are cater for uncertain expert knowledge in their representation and reasoning.
The  knowledge in these models is represented by causal relationships between variables and their probability measures.The alternate Bayesian approach uses the degree of a person's belief that an event will occur,rather than the actual probability that the event will occur.
\\\\\ The graphical representation of the dependencies embedded in the probabilistic models.The graphical representation used for belief propagation Techniques are two types.They are Markov  Random Fields and  Bayesian Networks.
Belief Propagation Techniques are used for performing inference on graphical models such as Bayesian Network and Markov Random Field.
 \\\\Belief Propagation Techniques are used in  wide variety of algorithms which are used in artificial intelligence,signal processing and digital communication.
\\\\ \textbf{ Keywords}:Belief Propagation Technique,Bayesian Network and Markov Random Field
\end{abstract}

%--------------------------------------------------------------------%
% CONTENTS, TABLES, FIGURES
\tableofcontents
\listoffigures

%--------------------------------------------------------------------%
%  Single counter for theorems and theorem-like environments:
\newtheorem{theorem}{Theorem}[chapter]
\newtheorem{assertion}[theorem]{Assertion}
\newtheorem{claim}[theorem]{Claim}
\newtheorem{conjecture}[theorem]{Conjecture}
\newtheorem{corollary}[theorem]{Corollary}
\newtheorem{definition}[theorem]{Definition}
\newtheorem{example}[theorem]{Example}
\newtheorem{figger}[theorem]{Figure}
\newtheorem{lemma}[theorem]{Lemma}
\newtheorem{prop}[theorem]{Proposition}
\newtheorem{remark}[theorem]{Remark}

%--------------------------------------------------------------------%
% Make the page numbers Arabic (1, 2, etc)
\cleardoublepage%
\pagenumbering{arabic}
